\chapter{Introduction}
\label{chap:introduction}

An absence seizure is short. It lasts only a few seconds, it shows up across the whole scalp
at once rather than in one spot, and on a long EEG recording it is genuinely easy for a human
reader to miss. Clinically, it is described as a generalized, roughly 3~Hz spike-and-wave
discharge, usually lasting 3 to 20 seconds, and it is seen mostly in children. Catching every
one of these events by eye, in hours of recording, is tedious and error-prone work.

This project builds an automatic detector for absence seizures, trained and evaluated on 21
real clinical EEG recordings. The goal was never simply to produce a high score on a benchmark;
it was to build something a clinician could actually trust: a detector that says, for a given
recording, here is where I believe a seizure happened, and here is how often I am wrong. That
second half --- being honest about false alarms, not just about seizures found --- shaped
almost every decision in this project, from how the data was split to how a detection is
scored.

\section{Objectives}

Three things were set out from the start:

\begin{enumerate}[leftmargin=1.4em]
  \item \textbf{Build a reliable classical pipeline.} Take the raw EEG recordings, clean and
    align them against their seizure annotations, filter and window the signal, extract
    hand-crafted features, and train a set of classical machine-learning models to detect
    seizures event by event, not just window by window.
  \item \textbf{Push past the obvious baseline.} Once a working pipeline exists, try to improve
    it: rebalancing the training data, augmenting the rare seizure class, and blending several
    models together, and report honestly which of these actually helped.
  \item \textbf{Compare against deep learning.} Train a convolutional neural network directly on
    the raw filtered signal, with no hand-crafted features at all, and measure it against the
    classical pipeline under the same evaluation rules, to see whether a model that learns its
    own representation can do better than one built on 320 hand-designed numbers per window.
  \item \textbf{Make the results usable.} Package everything behind a web application so that
    the recordings, the model predictions, and the comparison between models can actually be
    explored, not just read off a table in a report.
\end{enumerate}

\section{What this report covers}

The chapters that follow tell this story in the order it actually happened. \Cref{chap:rawdata}
covers the raw data itself: 22 recordings that started as raw MATLAB files with no declared
sampling rate, no declared amplitude unit, and seizure times stored as clock times with no
absolute start time to anchor them to --- and the concrete decisions made to turn that into a
clean, trustworthy 21-recording, 93-seizure dataset. \Cref{chap:preprocessing} covers the signal
processing and feature extraction built on top of that clean data.
\Cref{chap:classicalml,chap:augmentation,chap:ensemble} cover the classical machine-learning
pipeline: twelve models, three training strategies, and what helped and what did not.
\Cref{chap:deeplearning} covers the deep-learning comparison, including a real bug that produced
a real crash and how it was found and fixed. \Cref{chap:evaluation} pulls the evaluation
methodology together across all of the above, since several different threshold-selection
protocols are used in this project and keeping them straight matters for a fair comparison.
\Cref{chap:application} describes the web application built to explore the results.
\Cref{chap:lessons} is an honest account of what went wrong along the way, and
\Cref{chap:conclusion} closes with what the project actually achieved and what is still left to
do.

\begin{summarybox}
This project detects absence seizures in EEG, event by event, starting from 22 raw clinical
recordings with no declared sampling rate or amplitude unit and an annotation workbook that
needed real detective work to align. The result is a classical machine-learning pipeline (12
models), a deep-learning comparison, and a web application to explore both --- reported here
with every number traced back to its source, and every negative result kept in, not hidden.
\end{summarybox}
